% GENERATED FILE. Edit canonical lessons, then run npm run generate:books.
% canonical-source-hash: fnv1a64:29520b612c4d2f3b
% canonical-lessons: ES-C450-patio, ES-C450-hierba, ES-C450-piedra, ES-C450-cesped, ES-C450-entrenamiento, ES-R450-repaso-suelo, ES-C450-sintesis-suelo

\chapter{El césped y el patio}
\label{ch:es-suelo}
\hlchaptermodality{450}

\begin{chapteropening}
\textbf{By the end of this chapter:} \emph{I can read a notice pinned to the gate of a sports ground: say where training has moved to, tell a maintained pitch from grass that is merely growing, and warn somebody about stones on a path.}

Read the message or notice this chapter's words exist for, then say the same thing back.
\end{chapteropening}

\section[el patio]{el patio — the space in the middle}

\label{lesson:ES-C450-patio}



\begin{quote}
Say \textbf{el jardín} and \textbf{la calle}. A patio is neither, and it sits between them.
\end{quote}

\subsection*{You'll want to know: el patio}
\textbf{el patio} — the courtyard: an open space with the building around it.

\emph{Un piso con patio} — a flat with a courtyard. \emph{Los niños están en el patio} — the children are out in the yard. \emph{El patio de butacas} — the stalls of a theatre, the flat part with the seats.

A \emph{jardín} has plants and can be anywhere. A \emph{patio} is defined by what surrounds it, and it may hold nothing but a washing line.

\begin{grammarlens}[title={Grammar Lens: the school patio}]
In a school, \emph{el patio} is the playground, and with it comes a phrase you will hear constantly:

\noindent
\begin{tabularx}{\linewidth}{@{}>{\raggedright\arraybackslash}X>{\raggedright\arraybackslash}X@{}}
\toprule
\textbf{Spanish} & \textbf{what it covers} \\
\midrule
\emph{salir al patio} & to go out at break \\
\emph{la hora del patio} & break time \\
\bottomrule
\end{tabularx}

The second one names the time by naming the place, and Spanish does this often — \emph{la hora de la calle} for knocking-off time works the same way. When a Spanish phrase seems to be about a location, check whether it is really about a moment in the day.
\end{grammarlens}

\begin{cousinweb}
Here the honest answer is that nobody knows.

The best-supported account brings it from Occitan \emph{pati}, waste ground, and from there back to Latin \textbf{pactum}, an agreement — common land held by agreement. A rival account points at Latin \emph{patere}, to lie open, which fits the meaning better and the sound worse. Neither has settled the matter.

That is why this lesson lists no root. An empty list is not an oversight in this material; it is a claim, and the claim is that the origin is genuinely unresolved. Saying so is more useful than picking the prettier story, because the next time you see a root listed you can trust that somebody checked.

What is certain is where the word went. Spanish carried \emph{patio} into English, French and German, and it is now the ordinary word for a paved space behind a house in places that have never used it for anything else.
\end{cousinweb}

\subsection*{Your turn}
Say these aloud:

\begin{itemize}
\raggedright
  \item a flat with a courtyard
  \item the children are out at break
  \item the difference between a \emph{patio} and a \emph{jardín}
\end{itemize}

\begin{tcolorbox}[breakable,colback=teal!4,colframe=teal!35!black,title={Before you move on}]
The courtyard? (\textbf{El patio}.) What defines one? (\textbf{The building around it}.) What does \emph{la hora del patio} mean? (\textbf{Break time}.)
\end{tcolorbox}
\section[la hierba]{la hierba — grass, the plant}

\label{lesson:ES-C450-hierba}



\begin{quote}
Say \textbf{la planta} and \textbf{verde}. Now name the smallest, commonest green plant there is.
\end{quote}

\subsection*{You'll want to know: la hierba}
\textbf{la hierba} — grass, and more widely any low soft plant: a herb, a weed.

\emph{Hay hierba entre las piedras} — there is grass growing between the stones. \emph{Las malas hierbas} — the weeds, literally the bad grasses. \emph{Una infusión de hierbas} — a herbal tea.

You will also see it written \textbf{yerba}, mostly in the Americas, and \emph{yerba mate} is the usual spelling for the drink. Both are the same word said the same way.

\begin{grammarlens}[title={Grammar Lens: the h that is not pronounced and not decorative}]
Latin \emph{herba} had an h that had already gone silent by Roman times. Spanish kept the letter and added it back where Latin had none:

\noindent
\begin{tabularx}{\linewidth}{@{}>{\raggedright\arraybackslash}X>{\raggedright\arraybackslash}X>{\raggedright\arraybackslash}X@{}}
\toprule
\textbf{Latin} & \textbf{Spanish} & \textbf{the h} \\
\midrule
herba & la \textbf{h}ierba & inherited \\
ovum & el \textbf{h}uevo & added \\
ossum & el \textbf{h}ueso & added \\
\bottomrule
\end{tabularx}

Look at what the added ones have in common: they all begin \textbf{ue}. Scribes put an \emph{h} in front because a word starting with \emph{u} was read as a \emph{v}, and \emph{uevo} would have been read \emph{vevo}. The \emph{h} is a fence around the vowel.

\emph{Hierba} begins \emph{ie} rather than \emph{ue}, so it did not need the fence. It kept what Latin had already handed it. The letter is silent in all of them.
\end{grammarlens}

\begin{cousinweb}
Latin \textbf{herba}, grass or herb. English borrowed the same word twice: once through French as \textbf{herb}, and once as the botanical \textbf{herbaceous}. Spanish has the learned form too, in \emph{herbicida} and \emph{herbívoro}, where the h sits in front of an unbroken \emph{e} because those words came from books rather than from speech.

That is the tell. \emph{Hierba} has the broken vowel because ordinary speakers wore it down over centuries; \emph{herbívoro} has not, because it was lifted from Latin whole and handed to a scientist.
\end{cousinweb}

\subsection*{Your turn}
Say these aloud:

\begin{itemize}
\raggedright
  \item there is grass growing in the courtyard
  \item the weeds, using the Spanish phrase
  \item why \emph{huevo} has an h and \emph{hierba} has a different reason
\end{itemize}

\begin{tcolorbox}[breakable,colback=teal!4,colframe=teal!35!black,title={Before you move on}]
Grass, the plant? (\textbf{La hierba}.) Is the h pronounced? (\textbf{No} — not on its own, anywhere in Spanish.) What are \emph{las malas hierbas}? (\textbf{The weeds}.)
\end{tcolorbox}
\section[la piedra]{la piedra — the stone}

\label{lesson:ES-C450-piedra}



\begin{quote}
Say \textbf{la madera} and \textbf{duro}. Now name the building material that is harder than the first.
\end{quote}

\subsection*{You'll want to know: la piedra}
\textbf{la piedra} — the stone.

\emph{Una casa de piedra} — a stone house. \emph{Hay piedras en el camino} — there are stones on the path. \emph{De piedra} also works as a description of a person who has not moved: \emph{me quedé de piedra}, I was stunned.

For the material in the mass — quarried, cut, built with — Spanish stays singular: \emph{un muro de piedra}. For the things you can pick up, plural: \emph{piedras}.

\begin{grammarlens}[title={Grammar Lens: the same vowel trick, one row over}]
You have already watched a stressed \textbf{o} break into \textbf{ue} in \emph{rueda}. A stressed \textbf{e} does the same thing and becomes \textbf{ie}:

\noindent
\begin{tabularx}{\linewidth}{@{}>{\raggedright\arraybackslash}X>{\raggedright\arraybackslash}X>{\raggedright\arraybackslash}X@{}}
\toprule
\textbf{Latin} & \textbf{Spanish} & \textbf{what broke} \\
\midrule
petra & la p\textbf{ie}dra & e under stress \\
rota & la r\textbf{ue}da & o under stress \\
herba & la h\textbf{ie}rba & e under stress \\
\bottomrule
\end{tabularx}

Three of this chapter's words are the same sound change caught at different moments. Once you hear \emph{ie} and \emph{ue} as stressed vowels that gave way, a large part of Spanish spelling stops looking arbitrary.

And the test still works: \emph{petrificar} keeps the plain \textbf{e}, because the stress has moved off it.
\end{grammarlens}

\begin{cousinweb}
Latin \textbf{petra}, taken from Greek \textbf{pétra}, a rock.

Latin already had its own word for stone, \emph{lapis}, and it is the one that survived in the learned vocabulary — \emph{lapidario}, and English \emph{dilapidated}. The Greek borrowing won out in ordinary speech instead, which is the reverse of what usually happens.

The name \textbf{Pedro} is the same word: Greek \emph{pétros}, a stone, handed out as a nickname and then carried by enough people to become a name in its own right. So \emph{Pedro} and \emph{piedra} are one word twice, and a \textbf{petrol} is an oil that comes out of rock.
\end{cousinweb}

\subsection*{Your turn}
Say these aloud:

\begin{itemize}
\raggedright
  \item a stone house
  \item there is grass between the stones
  \item what \emph{piedra} and \emph{Pedro} have in common
\end{itemize}

\begin{tcolorbox}[breakable,colback=teal!4,colframe=teal!35!black,title={Before you move on}]
The stone? (\textbf{La piedra}.) Why \emph{ie} rather than \emph{e}? (\textbf{The stress fell on it}.) Which language did Latin take \emph{petra} from? (\textbf{Greek}.)
\end{tcolorbox}
\section[el césped]{el césped — grass you look after}

\label{lesson:ES-C450-cesped}



\begin{quote}
Say \textbf{el campo} and \textbf{plantar}. Somebody planted the grass on a football pitch, and that is the whole difference.
\end{quote}

\subsection*{You'll want to know: el césped}
\textbf{el césped} — the lawn, and on a sports ground the pitch.

\emph{Cortar el césped} — to mow the lawn. \emph{No pisar el césped} — keep off the grass, the wording on the sign in every park. \emph{El césped está seco} — the pitch is dry.

It has no everyday plural. Where English says \emph{the lawns}, Spanish repeats the singular or names the places instead.

\begin{grammarlens}[title={Grammar Lens: césped is not hierba}]
Both are grass. They answer different questions:

\noindent
\begin{tabularx}{\linewidth}{@{}>{\raggedright\arraybackslash}X>{\raggedright\arraybackslash}X@{}}
\toprule
\textbf{the question} & \textbf{the word} \\
\midrule
what plant is that? & \textbf{la hierba} \\
what surface is that? & \textbf{el césped} \\
\bottomrule
\end{tabularx}

\emph{Hierba} is the thing growing; \emph{césped} is the maintained expanse of it. So grass coming up between the paving stones is \textbf{hierba}, and nobody would call it \emph{césped}. A mown lawn is \textbf{césped}, though it is made of \emph{hierba}.

Sports tell the same story. \emph{Jugar en el césped} is to play on the pitch, and a pitch that has stopped being looked after has stopped being \emph{césped}.
\end{grammarlens}

\begin{cousinweb}
Latin \textbf{caespes}, a cut sod of turf — a slab of earth with grass on top, lifted with a spade for building banks and walls.

That is a good picture to hold, because it explains why \emph{césped} is about the surface rather than the plant. A \emph{caespes} was a thing you could carry.

The written accent is doing real work. A Spanish word ending in a consonant other than \emph{n} or \emph{s} takes its stress on the last syllable by default, so \emph{cesped} without a mark would be read \textbf{ces-PED}. The accent overrules the default and pins the stress on \textbf{CÉS}-ped.
\end{cousinweb}

\subsection*{Your turn}
Say these aloud:

\begin{itemize}
\raggedright
  \item I have to mow the lawn
  \item keep off the grass, as the sign puts it
  \item which word fits grass between the stones
\end{itemize}

\begin{tcolorbox}[breakable,colback=teal!4,colframe=teal!35!black,title={Before you move on}]
The lawn? (\textbf{El césped}.) Why the accent? (\textbf{Without it the stress would fall on the last syllable}.) A \emph{caespes} was? (\textbf{A cut sod of turf}.)
\end{tcolorbox}
\section[el entrenamiento]{el entrenamiento — the session}

\label{lesson:ES-C450-entrenamiento}



\begin{quote}
Say \textbf{el tren} and \textbf{correr}. One of those two is about to turn up somewhere you would not expect it.
\end{quote}

\subsection*{You'll want to know: el entrenamiento}
\textbf{el entrenamiento} — training, and a single training session.

\emph{Tengo entrenamiento los martes} — I have training on Tuesdays. \emph{El entrenamiento es en el césped de atrás} — training is on the back pitch. And the people: \textbf{el entrenador}, the coach; \textbf{entrenar}, to train; \textbf{entrenarse}, to train yourself.

Away from sport it keeps working: \emph{un entrenamiento de dos días} for a course at work.

\begin{grammarlens}[title={Grammar Lens: -miento names the process}]
You have met \textbf{-anza} turning a verb into a noun. \textbf{-miento} does the same job, and the two split the language between them:

\noindent
\begin{tabularx}{\linewidth}{@{}>{\raggedright\arraybackslash}X>{\raggedright\arraybackslash}X@{}}
\toprule
\textbf{verb} & \textbf{noun} \\
\midrule
entrenar & el entrena\textbf{miento} \\
mover & el movi\textbf{miento} \\
conocer & el conoci\textbf{miento} \\
\bottomrule
\end{tabularx}

Every \textbf{-miento} noun is masculine, as every \textbf{-anza} noun is feminine, so the ending hands you the article both times.

Which verb takes which is not predictable and has to be met word by word — \emph{esperanza} but \emph{movimiento}, with nothing in the verbs to tell you. What you can rely on is the gender, and that is worth having.
\end{grammarlens}

\begin{cousinweb}
French \textbf{entraîner}, to drag along — and inside it the same \emph{traîner} that gave Spanish \textbf{el tren}.

Hold the two together. A \textbf{tren} is the set of carriages \emph{dragged} behind the engine. To \textbf{entrenar} somebody is to \emph{drag them along} with you until they can keep up. The coach is out in front and the squad is the train.

English took the identical picture from the identical French word, which is why \emph{train} in English means both the thing on rails and what a coach does to a team. Spanish split the two forms — \emph{tren} and \emph{entrenar} — and so kept the relationship visible in a way English hides.
\end{cousinweb}

\subsection*{Your turn}
Say these aloud:

\begin{itemize}
\raggedright
  \item I have training on Tuesdays
  \item training is on the back pitch
  \item what a train and a coach are doing that is the same
\end{itemize}

\begin{tcolorbox}[breakable,colback=teal!4,colframe=teal!35!black,title={Before you move on}]
The training session? (\textbf{El entrenamiento}.) The coach? (\textbf{El entrenador}.) What gender do \emph{-miento} nouns take? (\textbf{Masculine}.)
\end{tcolorbox}
\section[(el patio, la hierba, la piedra, el …]{(el patio, la hierba, la piedra, el césped, el entrenamiento) — from cold}

\label{lesson:ES-R450-repaso-suelo}



\begin{quote}
Close the lessons before this one. Say all five: \emph{the courtyard}, \emph{grass}, \emph{the stone}, \emph{the lawn}, \emph{the training session}.
\end{quote}

\begin{grammarlens}[title={Grammar Lens: three words, one broken vowel}]
Say these aloud and listen for where the stress sits:

\noindent
\begin{tabularx}{\linewidth}{@{}>{\raggedright\arraybackslash}X>{\raggedright\arraybackslash}X@{}}
\toprule
\textbf{Latin or French} & \textbf{Spanish} \\
\midrule
herba & la h\textbf{ie}rba \\
petra & la p\textbf{ie}dra \\
entraîner & entren\textbf{a}miento \\
\bottomrule
\end{tabularx}

The first two broke their stressed \textbf{e} into \textbf{ie}. The third did not, because the stress in \emph{entrenamiento} never lands on that \textbf{e} at all — it lands three syllables later. Same vowel, same language, different outcome, and the only variable is stress.

This is why \emph{piedra} has a diphthong and \emph{petrificar} does not, and why \emph{rueda} has one and \emph{rodamos} does not. One condition explains all four.
\end{grammarlens}

\begin{grammarlens}[title={Grammar Lens: two pairs that are not synonyms}]
\noindent
\begin{tabularx}{\linewidth}{@{}>{\raggedright\arraybackslash}X>{\raggedright\arraybackslash}X@{}}
\toprule
\textbf{the pair} & \textbf{what separates them} \\
\midrule
patio / jardín & a patio is defined by what surrounds it \\
hierba / césped & hierba is the plant, césped is the tended surface \\
\bottomrule
\end{tabularx}

Both pairs fail in the same direction for an English speaker: \textquotedblleft{}yard\textquotedblright{} and \textquotedblleft{}grass\textquotedblright{} are wide words that cover each pair whole. Spanish is narrower here than English, which is the harder direction, because nothing in your own language warns you that a choice is waiting.
\end{grammarlens}

\subsection*{Your turn}
Say these aloud:

\begin{itemize}
\raggedright
  \item there are stones and grass in the courtyard
  \item somebody has to mow the lawn
  \item which of \emph{hierba} and \emph{césped} needs a gardener
\end{itemize}

\begin{tcolorbox}[breakable,colback=teal!4,colframe=teal!35!black,title={Before you move on}]
All five again, in Spanish, without looking. Then one sentence putting a training session on a pitch.
\end{tcolorbox}
\section[Practice]{(el aviso del polideportivo) — read it, then do it yourself}

\label{lesson:ES-C450-sintesis-suelo}



\begin{quote}
Say \textbf{el césped} and \textbf{el entrenamiento}. The notice below moves one of them off the other.
\end{quote}

\subsection*{Your turn}
A notice on the gate of the sports centre. Read it, then answer.

\begin{quote}
\textbf{AVISO}
\end{quote}

\begin{quote}
El \textbf{césped} del campo grande está en obras hasta el día 20.
\end{quote}

\begin{quote}
Los \textbf{entrenamientos} de esta semana se hacen en el \textbf{patio} de atrás.
\end{quote}

\begin{quote}
No dejen bolsas sobre la \textbf{hierba} del lateral. Se está replantando.
\end{quote}

\begin{quote}
Disculpen las molestias.
\end{quote}

Say these aloud:

\begin{itemize}
\raggedright
  \item where training happens this week, and why
  \item what you are asked not to do, and where
  \item when the big pitch will be available again
\end{itemize}

\begin{culture}
The notice uses \textbf{césped} once and \textbf{hierba} once, four lines apart, and swapping them would make it unreadable.

The big pitch is \textbf{césped}: a maintained playing surface, currently closed for work on it. The strip down the side is \textbf{hierba}: grass as a plant, being replanted. A Spanish reader takes \emph{se está replantando} as confirmation — you replant \emph{hierba}, since that is the living thing; \emph{césped} is what you get once the \emph{hierba} has taken and somebody starts mowing it.

Notice also \textbf{no dejen} and \textbf{disculpen}. Both are the \emph{ustedes} form, which is what public notices use: addressed to everyone, and formal without being cold. A sign saying \emph{no dejes} would read as though it had been written for one particular person who had already done it once.
\end{culture}

\begin{tcolorbox}[breakable,colback=teal!4,colframe=teal!35!black,title={Before you move on}]
Now your turn, out loud and then in writing. Four lines for a gate notice: say the pitch is closed, say where training moves to, ask people to keep off the grass at the side, and say there are stones on the path.

Then check yourself: did \emph{césped} go to the playing surface and \emph{hierba} to the plant, and did you use the \emph{ustedes} form?
\end{tcolorbox}
